\clearpage
\section{The three training recipes, from first principles}\label{sec:losses}
\subsection{What a training batch tells the model}
For a batch of $n$ normalized embeddings, form the similarity matrix
$S=ZZ^\top$, where row $i$ of $Z$ is $z_i^\top$. Let
\[
 y_{ij}=\ind[\text{action}_i=\text{action}_j].
\]
A \emph{positive pair} has the same action label; a \emph{negative pair} has
different labels. These names describe supervision, not good or bad examples.
The losses encourage a geometry consistent with those labels.

Each physical batch has eight action classes and four examples per class:
$P=8$, $K=4$, $n=PK=32$. Each anchor has three nonself positives and twenty-eight
negatives. The full matrix contains 96 directed positive pairs and 896 directed
negative pairs, excluding self-pairs. A directed pair counts $(i,j)$ separately
from $(j,i)$; the declared reductions handle that consistently.

Why not draw thirty-two arbitrary examples? An unbalanced random batch may have
too few same-label examples to provide the required positive relations. The
balanced sampler controls that issue and supplies the local population in which
contextual neighborhoods are defined.

\subsection{Standard contrastive: enforce similarity margins}
We want positive similarities to be sufficiently high and negative similarities
to be sufficiently low. Define a hinge $[a]_+=\max(a,0)$:
\[
 h^+_{ij}=[m_+-s_{ij}]_+,
 \qquad h^-_{ij}=[s_{ij}-m_-]_+.
\]
For the standard recipe, $m_+=0.9$ and $m_-=0.6$. A positive pair with
similarity $0.7$ receives a penalty $0.2$. A negative pair with similarity $0.7$
receives a penalty $0.1$. Positive similarity above $0.9$ and negative similarity
below $0.6$ produce no hinge penalty. The interval between the margins is an
intentional gap, not a probability threshold for an action classifier.

Define the strictly active sets
\[
 \mathcal A_+=\{(i,j):i\ne j,\ y_{ij}=1,\ h^+_{ij}>0\},\quad
 \mathcal A_-=\{(i,j):y_{ij}=0,\ h^-_{ij}>0\}.
\]
Then
\[
 L_{\mathrm{contrast}}=
 \frac{\sum_{(i,j)\in\mathcal A_+}h^+_{ij}}{|\mathcal A_+|}
 +\frac{\sum_{(i,j)\in\mathcal A_-}h^-_{ij}}{|\mathcal A_-|},
\]
where an empty active set contributes zero.

This detail is consequential: these are separate means over \emph{violations},
not one mean over all pairs. An abundance of already-satisfied negatives does
not dilute the remaining negative violations. The hinges are not squared.
Away from boundaries and changes in the active set, each active positive gives
derivative $-1/|\mathcal A_+|$ with respect to its similarity; each active
negative gives $+1/|\mathcal A_-|$. Through the chain rule these gradients change
$W,b$ so the relevant similarities move toward their margins.

\subsection{Context: ask how examples relate to the rest of the batch}
Direct similarity examines a pair. Contextual similarity also examines its
relations to other examples. Two noisy instances can be more convincingly
related if their neighborhoods agree. That is a motivation, not a guarantee:
bad neighborhoods can also reinforce bad relationships.

The current implementation has five steps: find nearby examples, measure shared
neighbors and non-neighbors, keep that agreement local, expand using reciprocal
neighbors, then make the score symmetric. The final score is trained to match
the same-action indicator $y_{ij}$. It is not trained to imitate the current
cosine matrix; doing that would change the target of the learning problem.

The following equations specify the local implementation. Their source trail
is the protocol and independently written local code; the contextual paper
provides the method's intellectual origin.\cite{contextual}

\subsubsection*{Step 1: define a neighborhood}
Compute squared unit-vector distances and a binary membership matrix:
\[
 d_{ij}=\max(2-2s_{ij},0),\qquad
 N_{ij}=\ind[d_{ij}\le d_{i,(k)}+\epsilon].
\]
$d_{i,(k)}$ is the $k$th smallest distance in row $i$. Here $k=4$ and
$\epsilon=0.05$. \textbf{Self is included} in the ordering. The tolerance
admits points just beyond the $k$th distance, so the neighborhood may contain
more than four examples, especially with ties.

The symbol $k$ is a neighbor rank used to set the distance cutoff; $K$ is
examples per class. Both equal
four in the fixed recipe, but their conceptual roles differ. Labels do not
choose the neighbors. Current embedding distances do.

\subsubsection*{Step 2: compare neighborhood membership}
For row $i$, set
\[
 a_i=\max\left(\sum_\ell N_{i\ell},1\right),\qquad
 b_i=\max\left(\sum_\ell(1-N_{i\ell}),1\right).
\]
The first-stage score is
\[
 A_{ij}=\frac12\left(
  \frac{\sum_\ell N_{i\ell}N_{j\ell}}{a_i}
 +\frac{\sum_\ell(1-N_{i\ell})(1-N_{j\ell})}{b_i}
 \right)N_{ij}.
\]
The numerators count shared neighbors and shared non-neighbors. The row
denominators turn them into fractions relative to $i$. Multiplication by
$N_{ij}$ restricts the score to pairs already in $i$'s neighborhood. The clamp
keeps an empty complement from creating division by zero; it is a declared
local finite convention.

\begin{concept}{Work one neighborhood calculation by hand}
In an invented six-example batch, let $i=1$ have neighbors $\{1,2,3,4\}$ and
$j=2$ have neighbors $\{1,2,3,5\}$. They share three of $i$'s four neighbors.
Their non-neighbor sets are $\{5,6\}$ and $\{4,6\}$, sharing one of two.
Because $N_{12}=1$,
\[
 A_{12}=\tfrac12(3/4+1/2)=0.625.
\]
This is only the first-stage score, not the final contextual similarity or
the cosine similarity. With unequal row sizes, $A_{ij}$ need not equal $A_{ji}$.
\end{concept}

\subsubsection*{Step 3: reciprocal expansion and symmetry}
Construct another neighborhood $N^{1/2}$ using the $k/2$ threshold and the
same tolerance. A reciprocal relation requires both directions:
\[
 R_{i\ell}=N^{1/2}_{i\ell}N^{1/2}_{\ell i}.
\]
Use those reciprocal neighbors to average first-stage rows, then symmetrize:
\[
 B_{ij}=\frac{\sum_\ell R_{i\ell}A_{\ell j}}
 {\max(\sum_\ell R_{i\ell},1)},\qquad
 C_{ij}=\frac{B_{ij}+B_{ji}}2.
\]
This incorporates the local relationship profiles of mutually close examples.
Symmetry gives a common score for the pair regardless of direction.

\subsubsection*{Step 4: fit action-label agreement}
The contextual penalty is
\[
 L_{\mathrm{context}}=\frac1{n^2}\sum_{i\ne j}(C_{ij}-y_{ij})^2.
\]
It asks contextual similarity to approach one for same-action pairs and zero
for different-action pairs. Self terms are omitted, but the denominator remains
$n^2$, not $n(n-1)$. Even a small reduction change can alter gradient scale and
therefore the relative weight of the objectives.

\subsection{How can a hard neighborhood comparison have a gradient?}
An indicator is locally constant almost everywhere and discontinuous at its
boundary. Its ordinary derivative cannot provide useful smooth updates for
moving points between neighborhoods. The recipe deliberately uses a
\textbf{straight-through estimator} (STE): perform the hard comparison in the
forward pass, but prescribe a surrogate derivative in the backward pass.

For the comparison $H(x,y)=\ind[x\ge y]$, the implemented rule assigns
\[
 \frac{\widetilde{\partial}H}{\partial x}=\alpha,\qquad
 \frac{\widetilde{\partial}H}{\partial y}=-\alpha,\qquad \alpha=10.
\]
The tilde emphasizes that this is a chosen backward rule, not the ordinary
derivative of the hard function. In the neighborhood expression, $x=-d_{ij}
+\epsilon$ and $y=-d_{i,(k)}$.

The $k$th-distance thresholds and first-stage counts are \emph{detached}: they
retain their forward values but are treated as constants during differentiation.
The reciprocal denominator in $B$ is \emph{not} detached. It participates in the
backward calculation. These choices are part of the fixed method, not details
that an optimizer can infer from a prose description.

\begin{caution}{What a derivative check can and cannot establish}
For a smooth head calculation, finite differences can check an analytic
gradient. For this hard comparison, a finite difference will usually see zero
while the prescribed STE returns a nonzero value. That disagreement is expected.
Verification must check the declared surrogate backward rule independently;
it must not claim the rule is the true derivative of the hard forward function.
\end{caution}

\subsection{The full contextual recipe includes two more terms}
The recipe also contains the margin contrastive loss and a mean-similarity
regularizer:
\[
 L_{\mathrm{reg}}=
 \left(\frac1{n^2}\sum_{i,j}s_{ij}-0.25\right)^2,
 \qquad
 L_{\mathrm{full}}=0.4L_{\mathrm{context}}
                  +0.6L_{\mathrm{contrast}}+0.1L_{\mathrm{reg}}.
\]
The regularizer includes diagonal entries. Its target controls average batch
similarity; it is not an action probability or a guarantee against collapse.
The contrastive term here uses $m_+=0.75$ and $m_-=0.6$, whereas the standalone
contrastive arm uses $m_+=0.9$.

That means a contextual win could arise from the combined effects of neighborhood
structure, a different margin, and regularization. Separating their contributions
would require matched ablations, including the omitted control. This study
instead compares the two complete declared recipes within its available scope.

The fixed values transfer a published recipe into a pose study. They were not
established as optimal pose hyperparameters. Do not substitute defaults from
an older teaching class or historical protocol for the explicit v3 configuration.

\subsection{SupCon: rank all positive examples against the batch}
Let $P(i)=\{p\ne i:y_{ip}=1\}$. SupCon assigns a softmax distribution over
every nonself example and penalizes the log probability of each positive:
\[
 L_{\mathrm{SupCon}}=-\frac1n\sum_i\frac1{|P(i)|}
   \sum_{p\in P(i)}\log
   \frac{\exp(s_{ip}/\tau)}{\sum_{a\ne i}\exp(s_{ia}/\tau)},
 \qquad \tau=0.07.
\]
For one anchor, every positive should rank well relative to all alternatives.
The denominator includes other positives as well as negatives; self is excluded.
Using the log converts a product-like probability objective into a sum and
strongly penalizes a positive assigned very small probability.

Temperature $\tau$ sets sensitivity to score gaps. Two alternatives whose
similarities differ by $0.14$ have softmax weight ratio
$\exp(0.14/0.07)=\exp(2)\approx7.39$. A smaller temperature makes the same
gap matter more. This does not make the value a calibrated probability that
an arbitrary final query has a particular action.

The study uses a single-view adaptation: no second augmented view and no extra
temperature multiplier. SupCon is supervised because action labels define
positives. It is not being used as a label-free pretraining method here.
In the ideal limit of zero probability on negatives, equal allocation among
three positives gives per-anchor loss $\log 3$ rather than zero. Raw loss values from different
recipes are therefore not a common accuracy scale.\cite{supcon}

\subsection{What fairness does and does not require}
All recipes receive the same architecture, data, batch plans, initialization
within a seed, training duration, and learning-rate grid. They are allowed to
use different objective formulas because that is the treatment being compared.
Equal tuning opportunities do not mean equal floating-point operation counts,
nor that each method reaches its own global optimum.

\begin{tabularx}{\linewidth}{@{}Y Y Y@{}}
\toprule Recipe & Main training signal & Interpretation boundary\\\midrule
Contrastive & Fix positive/negative margin violations & A specified baseline,
not every possible contrastive method.\\
Full contextual & Neighborhood agreement, contrastive margins, mean similarity
& Effects belong to the full recipe.\\
SupCon & Positive log probabilities against the batch & Descriptive reference;
no additional superiority test.\\\bottomrule
\end{tabularx}

\studycheck{Why is gradient accumulation from four batches of eight examples
not equivalent to one physical batch of thirty-two? Why would a contextual
win not isolate the neighborhood term? Why should loss magnitudes not be ranked
across recipes? Answers: Questions 5--7.}
\source{\src{src/pose_embed/losses/pairwise.py}{Contrastive and SupCon formulas};
\src{src/pose_embed/losses/contextual.py}{Local contextual implementation};
\src{docs/protocol/contextual-v3-source-notes.md}{Equation and recipe decisions};
\src{configs/protocol.v3.yaml}{Locked v3 settings}.}
